\documentclass[11pt]{article}
\usepackage[T1]{fontenc}
\usepackage{libertine}
\usepackage[margin=1in]{geometry}
\usepackage{amsmath,amssymb,amsthm,mathtools}
\usepackage[libertine]{newtxmath}
\usepackage[scaled=0.96]{zi4}
\usepackage{microtype,enumitem,etoolbox,xcolor,tikz,listings,tcolorbox,aliascnt}
\usepackage[hidelinks]{hyperref}
\usepackage[capitalise,nameinlink,noabbrev]{cleveref}

% Lean boxes, as in edemaine/translation-tiling-paper (paper.texlish, lines 50--68), with a
% text label in place of the Lean logo.
\definecolor{figblue}{RGB}{41,110,166}
\definecolor{figgreen}{RGB}{38,128,89}
\lstdefinelanguage{lean}{
  morekeywords={import,public,abbrev,def,theorem,inductive,structure,where,variable,
    match,with,fun,let,if,then,else},
  morekeywords=[2]{Prop,Type,Set,List,Fin},
  sensitive=true, morecomment=[l]{--}, morestring=[b]"}
% The literate list must cover LITERATE in extract_lean_snippets.py. Each symbol takes one
% column, so that comments aligned in the Lean source stay aligned; arrows take two (no
% aligned line has one).
\newcommand{\leansym}[1]{\makebox[0.5em]{$#1$}}
\newcommand{\leanwide}[1]{\makebox[1em]{$#1$}}
\lstdefinestyle{lean}{language=lean, basicstyle=\ttfamily\fontsize{9.5}{11.5}\selectfont, columns=fixed, basewidth=0.5em,
  keepspaces=true, showstringspaces=false, extendedchars=false,
  keywordstyle=\bfseries\color{figblue}, keywordstyle=[2]\color{figgreen},
  commentstyle=\upshape\color{black!60},
  literate={ℕ}{{\leansym{\mathbb{N}}}}1 {→}{{\leanwide{\to}}}2 {∀}{{\leansym{\forall}}}1 {∃}{{\leansym{\exists}}}1
    {∈}{{\leansym{\in}}}1 {∉}{{\leansym{\notin}}}1 {∧}{{\leansym{\wedge}}}1 {≤}{{\leansym{\le}}}1
    {↔}{{\leanwide{\leftrightarrow}}}2 {⟨}{{\leansym{\langle}}}1 {⟩}{{\leansym{\rangle}}}1 {α}{{\leansym{\alpha}}}1
    {`}{{\char0}}1}  % the T1 grave accent: zi4 has no \textasciigrave
\tcbuselibrary{skins}
\newcommand{\leanlabel}{\textsf{\bfseries\color{figblue}Lean}}
\newtcolorbox{leanbox}{colback=black!3,colframe=black!45,boxrule=0.5pt,arc=2pt,
  left=6pt,right=6pt,top=3pt,bottom=3pt,enhanced,
  overlay={\node[anchor=north east,inner sep=2pt,draw=black!45,fill=white,line width=0.5pt,
    rounded corners=2pt,font=\scriptsize] at (frame.north east) {\leanlabel};}}
\newcommand{\leanitem}[1]{\leansep\lstinputlisting[style=lean]{lean-snippets/#1.lean}\def\leansep{\vspace{2pt}}}
\newcommand{\leansnippet}[1]{\begin{leanbox}\def\leansep{}\let\do\leanitem\docsvlist{#1}\end{leanbox}}

\newcommand{\N}{\mathbb N}
\newcommand{\lean}[1]{\texttt{#1}}
\newcommand{\file}[1]{\texttt{#1}}
\setlength{\emergencystretch}{2em}
\setlist{itemsep=3pt,topsep=5pt}

\title{The Word RAM in Lean}
\author{}
\date{}

\begin{document}
\maketitle

\section{The Word RAM in Lean}\label{sec:lean}

This document explains the word RAM defined in the Lean file \file{Abstract.lean}, so that a
reader can check its definitions.
The boxes quote it and \file{AbstractBridge.lean} at commit \texttt{939256e6}; both files are
in \file{Complexity/Computability/WordRAM/}.

\subsection{Imports}

\file{Abstract.lean} imports only Mathlib and CSLib's base module \lean{Cslib.Init}, and
defines the machine from scratch.
(The source file remains the authoritative version.)

\leansnippet{Imports}

\subsection{Words and Operations}

At word size $w$, a \emph{word} is a natural number below $2^w$.
The machine stores arbitrary natural numbers, but an instruction faults when it would
create a value that is not a word (\cref{sec:lean-step}).

There are ten binary operations.
Subtraction is truncated ($3-5=0$).
Division and remainder round down, with Lean's conventions $x/0=0$ and $x \bmod 0 = x$.
The operations \lean{and}, \lean{or} and \lean{xor} act bitwise,
and \lean{shl} and \lean{shr} shift $x$ left and right by $y$ bits,
that is, give $x\cdot 2^y$ and $\lfloor x/2^y\rfloor$.
\lean{AluOp.eval} gives the value of an operation on natural numbers, before any check that
it is a word.

\leansnippet{AluOp,AluOp.eval}

\subsection{Instructions and Programs}

A machine has $r$ registers $R[0],\dots,R[r-1]$ and a memory cell $M[a]$ at every address
$a\in\N$.
The number $r$ of registers is a parameter of the types below (the implicit variable
\lean{r}).
There are ten instructions; the comment on each gives its effect.
Register arguments $d,i,j,a$ range over \lean{Fin r}${}=\{0,\dots,r-1\}$ (in \lean{load}
and \lean{store}, register $a$ holds the address); a constant $x$ and a jump target $t$ are
arbitrary natural numbers.
The instruction \lean{width} reads the word size $w$, so a program can depend on it.
A \emph{program} is a finite list of instructions, numbered from $0$.

\leansnippet{variable,Instr,Program}

\subsection{States and One Step}\label{sec:lean-step}

A \emph{state} consists of the number \lean{pc} of the next instruction, the registers and
the memory.
The memory is a function on all of $\N$: there is no allocation and no bound on addresses.

One step of a program continues in a new state, halts with an output, or faults.
It executes instruction number \lean{pc}:
\begin{itemize}
  \item \lean{const} and \lean{alu} write their value to $R[d]$ if it is a word, and fault
    otherwise. There is no wrap-around modulo $2^w$.
  \item \lean{width} writes $w$, which is always a word, since $w<2^w$.
  \item \lean{copy}, \lean{load} and \lean{store} move a value without a check.
  \item \lean{jump} and the two branches set \lean{pc}; every other instruction except
    \lean{halt} moves on to instruction $\lean{pc}+1$.
  \item \lean{halt}~$i$ stops with output $R[i]$.
  \item If \lean{pc} is not the number of an instruction, that is, $\lean{pc}$ is at least
    the length of the program, the step faults.
\end{itemize}
Copies, loads and stores need no check: when the input fits (\cref{sec:lean-input}), an
induction along the run shows that every register and memory cell always holds a word, and
so every address used is below $2^w$.
(This induction is not stated in the abstract files; it only explains why the check is not
needed. The library states it for its own word RAM as \lean{Cslib.WordRAM.Runs.bounded\_init}.)

\leansnippet{State,Outcome,step}

\subsection{Running Time}

\lean{Halts w p s v n} says that from state $s$, at word size $w$, program $p$ halts with
output $v$ after exactly $n$ instructions, the final \lean{halt} included.
It follows \lean{step} one instruction at a time: a fault ends the run without halting, and
a run that never halts satisfies \lean{Halts} for no $v$ and $n$.
Every instruction takes one unit of time, whatever its operands and the word size, so $n$ is
the running time.

\leansnippet{Halts}

\subsection{Input}\label{sec:lean-input}

An input is a list $x_1,\dots,x_k$ of natural numbers.
The start state is at instruction $0$, has all registers $0$, and has $k,x_1,\dots,x_k$ in
memory at addresses $0,1,\dots,k$ and $0$ at every other address;
so a program reads the length of its input at $M[0]$.
The input \emph{fits} word size $w$ if its length $k$ and all its entries are words.
The output of a run is the single number $R[i]$ at its \lean{halt}~$i$.

\leansnippet{init,InputFits}

\subsection{Decision in Time \texorpdfstring{$T$}{T}}

A decision problem is a set $A\subseteq\alpha$ of instances, given to the machine by an
encoding $\mathit{encode}\colon\alpha\to\text{\lean{List}}~\N$.
Two functions on instances set the terms: $T(a)$ bounds the running time on instance $a$,
and $\mathit{bits}(a)$ is the word-size threshold: the program must work at every word size
from $\mathit{bits}(a)$ on that the input fits.
A program $p$ \emph{decides $A$ in time $T$} if, for every instance $a$ and every word size
$w\ge \mathit{bits}(a)$ that $\mathit{encode}(a)$ fits, $p$ halts within $T(a)$ instructions,
with output $1$ if $a\in A$ and $0$ if not.
One program serves every such word size, and can read it with \lean{width}.
Word sizes that the input does not fit impose nothing, so the input need not fit at
$w=\mathit{bits}(a)$ itself.
$A$ is \emph{decidable in time $T$} if some program, with some number of registers, decides it.
(The length condition gives $k<2^w$, so all input addresses $0,\dots,k$ fit in a word; for
$k>0$ this is the same as $w>\log_2 k$. \lean{bits} can ask for more.)

\leansnippet{Decides,DecidableInTime}

\subsection{Main Theorems}\label{sec:lean-main}

Two theorems say that the machine of this section is the library's word RAM
\lean{Cslib.WordRAM} over its standard basis.
They are stated in \file{AbstractBridge.lean}, which imports both \file{Abstract.lean} and the
library's word RAM, in the namespace \lean{Cslib.WordRAM.AbstractBridge}.
There, names with the prefix \lean{Abstract.} are those of this section,
and names without it are from the library's word RAM \lean{Cslib.WordRAM}:
\lean{Program .standard r} is a program over the standard basis with $r$ registers,
a state \lean{$\langle$pc, $\langle$reg, mem$\rangle\rangle$} has the same three parts as here,
\lean{step} is one step,
and \lean{Executes w q s v n} says that $q$ halts from $s$ with output $v$ after exactly $n$
instructions.
A reader does not need to know how these library definitions are made:
\lean{decidableInTime\_iff} makes the library's decidability in time $T$ equal to that of
this section, and \lean{runsAgree} does the same for single steps and halting runs.

Two outcomes, one of each machine, \emph{agree} if both continue in the same state, both halt
with the same output, or both fault.
An abstract program $p$ and a library program $q$ \emph{run alike} if, from every state and at
every word size, one step of each gives agreeing outcomes, and $p$ halts with output $v$
after $n$ instructions exactly when $q$ does.

\leansnippet{OutcomeAgrees,Same}

For every number $r$ of registers, every abstract program runs alike with some library
program, and every library program runs alike with some abstract program.
(In the proof, the program of the other machine is a translation instruction by
instruction.)

\leansnippet{runsAgree}

The two machines have the same time classes: for every type $\alpha$ of instances, problem
$A$, encoding, time bound $T$ and word-size threshold $\mathit{bits}$, $A$ is decidable in time
$T$ on the machine of this section exactly when it is on the library's word RAM over the
standard basis.

\leansnippet{decidableInTime_iff}

So a library theorem that concludes
\lean{DecidableInTime .standard A encode T bits} is, by \lean{decidableInTime\_iff},
a time bound on the machine described in this section.

To trust these results, a reader needs to read only the definitions of \file{Abstract.lean}
and the statements of these two theorems, with \lean{OutcomeAgrees} and \lean{Same}.
Lean checks the proofs, which use only Lean's three standard axioms (\lean{propext},
\lean{Classical.choice}, \lean{Quot.sound}).

\end{document}
